
\SGASubsection{4.1}{Smooth, Connected, Affine Algebraic Groups}
\SGARefTarget{sga3:xvii:subsection:4:1}

\SGAStatement{\SGAInlineRefTarget{sga3:xvii:proposition:4-1-1}Proposition 4.1.1}{Let \(k\) be a field, let \(G\)
be a smooth, connected, affine \(k\)-algebraic group, and put
\(\mathfrak g=\operatorname{Lie}G\).  The following properties are
equivalent:
\begin{enumerate}
\item[(i)] \SGARefTarget{sga3:xvii:proposition:4:1:1:item:i}
  \(G\) is unipotent.
\item[(ii)] \SGARefTarget{sga3:xvii:proposition:4:1:1:item:ii}
  \(G\) has a central composition series whose successive quotients
  are forms of \(\mathbb G_a\).
\item[(iii)] \SGARefTarget{sga3:xvii:proposition:4:1:1:item:iii}
  \(G\) has a central, characteristic composition series whose
  successive quotients are forms of \((\mathbb G_a)^r\).
\item[(iv)] \SGARefTarget{sga3:xvii:proposition:4:1:1:item:iv}
  There is an integer \(n>1\) such that \(G_{\overline{k}}\)
  contains no subgroup isomorphic to \(\mu_n\).
\item[(v)] \SGARefTarget{sga3:xvii:proposition:4:1:1:item:v}
  Every maximal torus of \(G\) is the trivial group.
\end{enumerate}
Suppose, moreover, that \(G\) is an algebraic subgroup of a linear
group \(\operatorname{GL}(n)\).  The preceding conditions are then
also equivalent to:
\begin{enumerate}
\item[(vi)] \SGARefTarget{sga3:xvii:proposition:4:1:1:item:vi}
  \(\mathfrak g\subset\mathfrak{gl}(n)\) consists of nilpotent
  endomorphisms.
\item[(vii)] \SGARefTarget{sga3:xvii:proposition:4:1:1:item:vii}
  \(\mathfrak g\) is nilpotent, and its center contains no nonzero
  semisimple endomorphism.
\end{enumerate}}
\SGARefTarget{sga3:xvii:proposition:4:1:1}

\SGAProof{Proof}
The implication \SGARefLink{sga3:xvii:proposition:4:1:1:item:ii}{(ii)}\(\Rightarrow\)\SGARefLink{sga3:xvii:proposition:4:1:1:item:i}{(i)} is clear, and
\SGARefLink{sga3:xvii:proposition:4:1:1:item:i}{(i)}\(\Rightarrow\)\SGARefLink{sga3:xvii:proposition:4:1:1:item:iii}{(iii)} was proved in \SGARefLink{sga3:xvii:lemma:3-9}{3.9}.  The implication
\SGARefLink{sga3:xvii:lemma:3:9:item:iii}{(iii)}\(\Rightarrow\)\SGARefLink{sga3:xvii:lemma:3:9:item:ii}{(ii)} will follow from the next lemma.

\SGAStatement{Lemma 4.1.2}{Let \(k\) be a field, and let \(G\) be
a \(k\)-algebraic group that is a form of \((\mathbb G_a)^r\).
Then:
\begin{enumerate}
\item[(a)] \SGARefTarget{sga3:xvii:lemma:4:1:2:item:a}
  \(G\) can be realized as an algebraic subgroup of
  \((\mathbb G_a)^n\), for a suitable integer \(n\).
\item[(b)] \SGARefTarget{sga3:xvii:lemma:4:1:2:item:b}
  \(G\) has a composition series whose successive quotients are
  forms of \(\mathbb G_a\).
\end{enumerate}}
\SGARefTarget{sga3:xvii:lemma:4:1:2}

Indeed, by hypothesis there is an extension \(k'\) of \(k\) such
that \(G_{k'}\) is isomorphic to \((\mathbb G_{a,k'})^r\).  By the
principle of finite extension (EGA IV, 9.1.1), we may assume that
\(k'/k\) is finite.  For (a), it then suffices to consider the
canonical closed immersion (EGA V)
\[
  G\longrightarrow
  \prod_{k'/k}(G_{k'})_{/k'}
  \xrightarrow{\ \sim\ }(\mathbb G_{a,k})^n,
  \qquad n=r\deg(k'/k).
\]

To prove \SGARefLink{sga3:xvii:lemma:4:1:2:item:b}{(b)}, in view of \SGARefLink{sga3:xvii:lemma:4:1:2:item:a}{(a)}, we may assume that \(G\) is a closed
subgroup of \(G'=(\mathbb G_a)^n\).  If \(G\ne0\), there is a
hyperplane \(\mathfrak h\) in
\(\mathfrak g'=\operatorname{Lie}G'\) that does not contain
\(\mathfrak g\).  Let \(H\) be the vector subgroup
\(\mathbf W(\mathfrak h)\) of
\(\mathbf W(\mathfrak g')=G'\).  Since \(H\) is defined by one
equation in \(G'\), the intersection \(G\cap H\) is defined by one
equation in \(G\), and
\[
\begin{aligned}
  \dim G-1
  &\leq \dim(G\cap H)\\
  &\leq \dim\operatorname{Lie}(G\cap H)
   =\dim_k(\mathfrak g\cap\mathfrak h)\\
  &=\dim_k\mathfrak g-1
   =\dim G-1.
\end{aligned}
\]
Thus
\(\dim(G\cap H)=\dim\operatorname{Lie}(G\cap H)\), and consequently
\(G\cap H\) is smooth.  The group
\[
  G_1=(G\cap H)^0
\]
is a smooth, connected algebraic subgroup of \(G\) such that
\(G/G_1\) is smooth, connected, and one-dimensional, hence is a
form of \(\mathbb G_a\) (\SGARefLink{sga3:xvii:subsection:4:1}{4.1}, \SGARefLink{sga3:xvii:proposition:4:1:1:item:i}{(i)}\(\Rightarrow\)\SGARefLink{sga3:xvii:proposition:4:1:1:item:iii}{(iii)}).  Induction
on the dimension of \(G\) completes the proof.
